\subsection{UML Modelling of the Proposed DSS}
\label{sec:uml_modeling}

Before writing a single line of implementation code, we modelled the
proposed DSS using four Unified Modelling Language (UML) diagrams.
Each diagram answers a different question about the system: who uses
it, how it is structured internally, how its components communicate,
and in what order the work gets done.

%----------------------------------------------------------------
\subsubsection{Use Case Diagram}
%----------------------------------------------------------------

The use case diagram answers the question: \textit{who interacts with
the system, and what can each actor do?}  It defines the system
boundary and maps every major function to the role that is allowed
to trigger it.

The system defines five human roles organised in a strict privilege
hierarchy.  A \textbf{Viewer} has read-only access: they can monitor
live KPIs on the dashboard, consult the alert history, and view
optimisation summaries, but cannot initiate any action that modifies
system state.  An \textbf{Operator} can do everything a Viewer can,
and additionally request optimisation runs, interact with the AI
setpoint advisor, and acknowledge active alerts.  An
\textbf{Engineer} has full operational authority: they can review
optimisation recommendations, click the \textit{I Applied It} button
to confirm setpoint application, and consult the FOPDT process
tracking panel to verify that the plant responded as predicted.  A
\textbf{Company Administrator} controls the configuration scope of
their own plant: they can manage sites, distillation columns, model
parameters, and FOPDT dynamics settings.  A \textbf{System
Administrator} has unrestricted access, including user management
across companies and global application configuration.

In addition to human actors, four software actors interact with the
system.  The \textbf{DCS/SCADA data source} pushes real-time sensor
readings into the ingestion pipeline via WebSocket.  The
\textbf{ML prediction engine} evaluates the current process state
and returns simultaneous estimates of energy consumption and product
purity.  The \textbf{FOPDT simulation engine} is invoked after an
engineer confirms setpoint application: it computes the expected
plant trajectory at $+5$, $+15$, and $+30$ minutes using the
First-Order Plus Dead Time model for each controlled loop, and
stores the predicted trajectory for subsequent comparison with live
DCS readings.  The \textbf{rule-based monitoring module} checks
sensor health at each ingestion cycle and generates alerts.
Separating these as explicit actors makes clear that the DSS does
not operate in isolation --- it is embedded in a larger industrial
data ecosystem.

\begin{figure}[H]
    \centering
    \includegraphics[width=\textwidth]{figures/usecase.png}
    \caption{Use Case Diagram of the proposed DSS. The five human
             roles (Viewer, Operator, Engineer, Company Administrator,
             System Administrator) and four software actors
             (DCS/SCADA source, ML engine, FOPDT simulation engine,
             monitoring module) interact with seven functional groups:
             authentication, real-time monitoring, setpoint
             optimisation, post-apply FOPDT tracking, alert
             management, audit trail consultation, and system
             administration.}
    \label{fig:usecase}
\end{figure}

%----------------------------------------------------------------
\subsubsection{Class Diagram}
%----------------------------------------------------------------

The class diagram answers the question: \textit{what are the main
entities in the system, what data do they hold, and how are they
related?}  It is a static, technology-independent view of the
system's structure.

The model is organised around three layers: organisational entities,
process data entities, and analytical service classes.

\paragraph{Organisational entities.}
\textbf{Company} represents an industrial client, holding display
configuration (brand colours, logo) and a reference to the external
DCS endpoint from which sensor data is pulled.  \textbf{Site}
represents a physical plant or branch belonging to a company.
\textbf{DistillationColumn} represents one fractionation unit within
a site, with a unique tag identifier, a reference to its surrogate
model file, and a JSON configuration block that stores setpoint
bounds and sensor mappings.  This three-level hierarchy
(Company~$\rightarrow$~Site~$\rightarrow$~Column) allows the same
DSS instance to serve multiple industrial clients without any data
cross-contamination.  \textbf{User} stores credentials, role, and a
foreign-key link to the company they belong to.

\paragraph{Process data entities.}
\textbf{Prediction} stores the output of the ML engine for a given
ingestion cycle: predicted energy consumption, product purity, and
butane flow, together with an outlier flag and a stability score.
\textbf{OptimizationResult} stores the complete snapshot of each
optimisation run, including current and recommended setpoints, the
predicted energy savings and purity improvement, the identity of the
requesting user, and two application-tracking fields: an
\textit{applied} flag with an \textit{applied\_at} timestamp set
when the engineer clicks \textit{I Applied It}, and a
\textit{predicted\_trajectory} JSON column that persists the
per-horizon FOPDT simulation output at the moment of application so
that the tracking endpoint can compare it against live DCS readings
later.  \textbf{Alert} represents an anomaly event with a type,
severity level, and the sensor tag that triggered it.
\textbf{AuditLog} is an append-only record of every protected action
in the system, capturing the username, role, HTTP endpoint, action
label, and response code; this table is the technical foundation of
the full traceability feature.

\paragraph{Analytical service classes.}
\textbf{PredictionEngine} takes a process state and returns a
Prediction using the XGBoost surrogate model.
\textbf{OptimisationEngine} runs Differential Evolution over the
surrogate to return the recommended setpoints and gain estimates.
\textbf{FopdtSimulator} is invoked on apply: it uses the per-loop
$K$, $\tau$, and $\theta$ parameters configured by the administrator
to compute the expected tag values at each time horizon and stores
the result in the OptimizationResult.
\textbf{TrackingService} is called at each horizon after apply: it
fetches the stored trajectory, compares it with the live DCS
readings supplied by the frontend, computes per-tag deviation
percentages, and returns a tracking score and a re-optimise flag.
\textbf{AnomalyDetector} checks sensor health at each ingestion
cycle and generates Alerts.  Keeping the service classes separate
from the data entities means that any algorithm can be swapped ---
for example, replacing the XGBoost surrogate with a GRU --- without
changing the rest of the system.

\begin{figure}[H]
    \centering
    \includegraphics[width=\textwidth]{figures/classdia.png}
    \caption{Conceptual Class Diagram of the proposed DSS. The three
             organisational entities (Company, Site,
             DistillationColumn) anchor the multi-tenant structure.
             Process data entities (Prediction, OptimizationResult,
             Alert, AuditLog) record every system event. Analytical
             service classes (PredictionEngine, OptimisationEngine,
             FopdtSimulator, TrackingService, AnomalyDetector)
             encapsulate the algorithmic logic independently of the
             data layer.}
    \label{fig:classdiagram}
\end{figure}

%----------------------------------------------------------------
\subsubsection{Sequence Diagram}
%----------------------------------------------------------------

The sequence diagram answers the question: \textit{in what order do
the system components exchange messages during a real interaction?}
It shows the dynamic behaviour of the system across time, which the
class diagram cannot capture.

Five operational scenarios are represented in the diagram.  The
\textbf{authentication flow} shows the user submitting credentials,
the backend verifying them against the database and checking the
user's role and company membership, and a signed JWT token being
returned to the frontend and stored for subsequent requests.  The
\textbf{real-time data pipeline} shows the continuous ingestion
loop: at each cycle the backend receives a new sensor reading via
WebSocket, passes it to the prediction engine and the rule-based
monitor in parallel, stores the resulting Prediction in the
database, and broadcasts a live update to the frontend.  If the
monitor raises an alert, it is stored in the Alert table and an
AuditLog entry is created; the alert is pushed to the frontend in
the same broadcast cycle.  The \textbf{optimisation request flow}
shows the engineer requesting an optimisation from the dashboard:
the backend validates the JWT and role, calls the Differential
Evolution engine (which internally queries the XGBoost surrogate in
a loop), stores the result as an OptimizationResult, writes an
AuditLog entry, and returns the recommended setpoints and predicted
gains to the frontend.  The \textbf{apply and FOPDT tracking flow}
shows the engineer clicking \textit{I Applied It}: the backend marks
the OptimizationResult as applied, immediately invokes
FopdtSimulator to compute the predicted plant trajectory at $+5$,
$+15$, and $+30$ minutes, persists this trajectory in the
\textit{predicted\_trajectory} column, writes a second AuditLog
entry with the action label \texttt{APPLY\_RECOMMENDATION}, and
confirms success to the frontend.  The frontend then starts three
countdown timers; when each fires, it sends the current live DCS
tag values to the \texttt{/tracking} endpoint, which calls
TrackingService, compares them with the stored trajectory, and
returns per-tag deviations, a tracking score, and a
re-optimise flag that the frontend renders as a colour-coded status
panel.  The \textbf{alert acknowledgement flow} shows the operator
marking an alert as resolved, triggering a database update, an
AuditLog entry, and a confirmation back to the interface.

\begin{figure}[H]
    \centering
    \includegraphics[width=\textwidth]{figures/sequence.png}
    \caption{Sequence Diagram of the proposed DSS, covering five
             operational scenarios: user authentication with role
             verification, real-time data ingestion and parallel
             prediction, optimisation request and setpoint
             recommendation, post-apply FOPDT simulation and
             horizon-triggered process tracking, and alert
             acknowledgement. Every protected action produces an
             AuditLog entry.}
    \label{fig:sequencediagram}
\end{figure}

%----------------------------------------------------------------
\subsubsection{Activity Diagram}
%----------------------------------------------------------------

The activity diagram answers the question: \textit{what is the
overall operational workflow of the system, from startup to decision
output?}  Unlike the sequence diagram, which focuses on message
exchanges between components, the activity diagram emphasises
decision points and parallel flows within the system as a whole.

The workflow begins with user authentication and role resolution.  On
success, the system identifies the active column context from the
engineer's company and site assignment, enabling the multi-column
generalisation: if the company operates several distillation units,
the engineer selects the target column from a menu and all
subsequent model calls and setpoint recommendations are routed to
the column-specific surrogate and configuration.  The backend then
starts the ingestion service, which enters a continuous loop.  At
each iteration, a new process state is received from the DCS data
source; the prediction model and the rule-based anomaly detector run
in parallel on this state.  The prediction result is persisted to
the database and pushed to the frontend.  If the anomaly detector
finds a fault, the alert is classified by type and severity, stored,
logged to the AuditLog, and sent to the frontend in real time.

When the engineer requests a setpoint optimisation, the workflow
branches: the Differential Evolution engine generates candidate
setpoints, the XGBoost surrogate evaluates them, a feasibility check
discards candidates that violate the minimum purity or product-flow
constraints, and the iteration continues until the convergence
criterion is met or the maximum iteration count is reached.  The
best feasible setpoint is returned, stored as an OptimizationResult,
and displayed to the engineer, who reviews it before any change is
applied to the real process.

If the engineer confirms the recommendation by clicking \textit{I
Applied It}, the workflow enters the post-apply branch.  The FOPDT
simulator runs immediately, computing the expected tag values at
$+5$, $+15$, and $+30$ minutes for each controlled loop using the
First-Order Plus Dead Time step-response model; this trajectory is
stored alongside the OptimizationResult.  Three horizon timers are
started in the frontend.  At each horizon, the frontend queries the
live DCS readings and calls the tracking endpoint; the tracking
service compares the actual readings with the stored trajectory and
computes per-tag deviation percentages.  If the worst-case deviation
exceeds 40\% or the overall tracking score falls below 0.5, the
system raises a re-optimise flag, indicating a disturbance,
equipment fault, or FOPDT parameter mismatch that warrants a new
optimisation cycle.  Otherwise, the system confirms good tracking
and the loop continues until all three horizons have been evaluated.

\begin{figure}[H]
    \centering
    \includegraphics[width=\textwidth]{figures/activity.png}
    \caption{Activity Diagram of the proposed DSS, showing the
             end-to-end operational workflow from user login and
             column selection, through continuous data ingestion with
             parallel prediction and anomaly detection, constrained
             Differential Evolution optimisation, operator review,
             post-apply FOPDT simulation, and horizon-triggered
             process tracking with automatic re-optimise detection.}
    \label{fig:activitydiagram}
\end{figure}

%----------------------------------------------------------------
\subsubsection{Summary}
%----------------------------------------------------------------

The four diagrams provide complementary and non-overlapping views of
the proposed DSS.  The use case diagram establishes what the system
must do and for whom, capturing the five-role privilege hierarchy and
the new post-apply tracking and audit use cases introduced in this
work.  The class diagram defines what the system is made of,
highlighting the three-level multi-tenant organisational hierarchy
(Company, Site, DistillationColumn), the FOPDT trajectory
persistence in OptimizationResult, and the append-only AuditLog
that underpins full operational traceability.  The sequence diagram
shows how the components interact during execution, with particular
attention to the apply-and-track flow that distinguishes this system
from a simple recommendation tool.  The activity diagram describes
the full operational logic from startup to post-apply verification,
including the decision branch that triggers a new optimisation cycle
when the process deviates significantly from the predicted FOPDT
trajectory.  Together they ensure that the implementation phase
begins from a clearly specified and internally consistent design, and
that anyone reading this thesis can understand how the system works
without access to the source code.
